\section{Sección de carga y relaciones de cuantización eléctrica}
\label{iii:sec:carga}

La respuesta constitutiva establece una relación entre las dos hojas y una
escala relativa de propagación. La sección de carga se obtiene al componer
esta respuesta con la constante de estructura fina y la acción, ambas
construidas en los capítulos antecedentes. El orden tiene una consecuencia
precisa: la carga no determina retrospectivamente el registro dodecafásico.
Es el registro, junto con la acción y la respuesta del vacío, el que
determina la sección de carga en la recta correspondiente.

\subsection{Determinación positiva de la sección de carga}

Consideremos las secciones positivas $\hbar_a$, con
$a\in\{\mathrm{pre},\mathrm{ret}\}$, y $h_a=2\pi\hbar_a$.
La impedancia $Z_{\rm vac}$ ya está construida mediante
\eqref{iii:eq:dimensions}. El cociente $h_a/Z_{\rm vac}$ pertenece a
la recta cuadrática de carga. Esta compatibilidad dimensional permite
formular la siguiente definición sin elegir todavía una carta SI.

\begin{definition}[Sección positiva de carga]
La sección de carga correspondiente a la orientación de acción $a$ es
\begin{equation}
 e_a=\left(\frac{2\alpha h_a}{Z_{\rm vac}}\right)^{1/2}.
 \label{iii:eq:charge}
\end{equation}
La raíz se toma en la semirrecta positiva de carga. Su sección conjugada
es $-e_a$; el signo de esta última no se identifica con el índice de
retorno de la acción.
\end{definition}

\begin{proposition}[Existencia, unicidad y compatibilidad del acoplamiento]
La ecuación $Z_{\rm vac}e_a^2=2\alpha h_a$ determina una única sección
positiva. En la realización constitutiva se cumple también
\begin{equation}
 \alpha=\frac{Z_{\rm vac}e_a^2}{2h_a}
 =\frac{e_a^2}{4\pi\varepsilon_{\rm vac}\hbar_a c_{\rm vac}}.
 \label{iii:eq:coupling}
\end{equation}
\end{proposition}
\begin{proof}
La positividad de $\alpha$, $h_a$ y $Z_{\rm vac}$ garantiza una raíz
positiva única. Las identidades constitutivas dan
$\varepsilon_{\rm vac}c_{\rm vac}=Z_{\rm vac}^{-1}$.
Sustituir esta igualdad y $h_a=2\pi\hbar_a$ en la primera expresión
de \eqref{iii:eq:coupling} produce la segunda.
\end{proof}

La última expresión es la forma habitual del acoplamiento electromagnético
en una carta física. Aquí aparece como identidad posterior de las secciones
construidas. La prueba no vuelve a calcular $\alpha$ a partir de una carga
medida: establece la coherencia algebraica del resultado anterior con la
sección que acaba de definirse. La identificación física exige, además,
mantener la carta de unidades y el régimen del acoplamiento empleados en la
comparación. Una identidad entre magnitudes y una comparación experimental
tienen funciones diferentes en ese reconocimiento.

\subsection{Resistencia, conductancia, flujo y frecuencia}

Las siguientes cuatro magnitudes utilizan la misma sección $e_a$. Su
dependencia puede resolverse por completo antes de evaluar ningún decimal.

\begin{theorem}[Composición de los cuantos eléctricos]
Las secciones
\begin{align}
 R_K&=\frac{h_a}{e_a^2}=\frac{Z_{\rm vac}}{2\alpha},
 &G_0&=\frac{2e_a^2}{h_a}=\frac{4\alpha}{Z_{\rm vac}},
 \label{iii:eq:rk-g0}\\
 \Phi_{0,a}&=\frac{h_a}{2e_a}
 =\sqrt{\frac{h_aZ_{\rm vac}}{8\alpha}},
 &K_{J,a}&=\frac{2e_a}{h_a}=\Phi_{0,a}^{-1}
 \label{iii:eq:flux-josephson}
\end{align}
satisfacen exactamente
\begin{equation}
 R_KG_0=2,\qquad G_0Z_{\rm vac}=4\alpha,\qquad
 K_{J,a}\Phi_{0,a}=1,\qquad
 K_{J,a}^{2}R_K=\frac4{h_a}.
 \label{iii:eq:quanta-identities}
\end{equation}
\end{theorem}
\begin{proof}
La sustitución de $e_a^2=2\alpha h_a/Z_{\rm vac}$ en las dos primeras
definiciones cancela $h_a$ y da \eqref{iii:eq:rk-g0}. Al elevar al
cuadrado la sección positiva $h_a/(2e_a)$ se obtiene
$h_aZ_{\rm vac}/(8\alpha)$; la positividad fija la raíz en
\eqref{iii:eq:flux-josephson}. Su inversa es $2e_a/h_a$.
Los tres primeros productos de \eqref{iii:eq:quanta-identities} se
reducen, respectivamente, a $2$, $4\alpha$ y $1$. Finalmente,
$(4e_a^2/h_a^2)(h_a/e_a^2)=4/h_a$.
\end{proof}

La realización física reconoce en estas secciones la resistencia de von
Klitzing, el cuanto de conductancia, el cuanto de flujo y la constante de
Josephson. El factor $2$ de la conductancia pertenece a la normalización
de la magnitud aquí definida; su identificación con una multiplicidad de
canales de un sistema material exige especificar esa representación.
Las cuatro ecuaciones no constituyen cuatro generadores independientes:
comparten $\alpha$, la acción y la respuesta del vacío. Precisamente esa
dependencia produce las identidades cruzadas y permite comprobar una
evaluación sin alterar las operaciones que la originaron.

\subsection{Retorno de la acción y transporte de las secciones}

El cociente $R_{\rm act}=h_{\rm pre}/h_{\rm ret}$, obtenido antes del
par angular, conserva la diferencia entre las dos secciones de acción.
No cambia las asignaciones de hoja de la respuesta constitutiva. Son dos
operaciones distintas: retorno de la acción e intercambio de los canales.

\begin{proposition}[Covariancia bajo el retorno]
Con $\alpha$ y $Z_{\rm vac}$ comunes, se tiene
\begin{align}
 \frac{e_{\rm pre}}{e_{\rm ret}}&=R_{\rm act}^{1/2},
 &R_K^{\rm pre}&=R_K^{\rm ret},
 &G_0^{\rm pre}&=G_0^{\rm ret},\label{iii:eq:charge-return}\\
 \frac{\Phi_{0,\rm pre}}{\Phi_{0,\rm ret}}&=R_{\rm act}^{1/2},
 &\frac{K_{J,\rm pre}}{K_{J,\rm ret}}&=R_{\rm act}^{-1/2}.
 \label{iii:eq:flux-return}
\end{align}
\end{proposition}
\begin{proof}
La primera identidad sigue de dividir las dos versiones de
\eqref{iii:eq:charge} y tomar la raíz positiva. En
\eqref{iii:eq:rk-g0} la acción ya se ha cancelado. Las otras dos
identidades siguen de las potencias $h_a^{1/2}$ y $h_a^{-1/2}$ en
\eqref{iii:eq:flux-josephson}.
\end{proof}

Resistencia y conductancia conservan el cociente constitutivo, mientras
que flujo y frecuencia retienen también la orientación de acción. Esta
separación da una lectura estructural de sus dependencias: unas magnitudes
son invariantes del cociente de hojas y otras miden, además, la sección de
acción utilizada. El contraste metrológico debe declarar qué sección evalúa;
la proximidad numérica de las dos acciones no autoriza identificarlas.
